\documentclass[12pt]{article}
\usepackage[utf8]{inputenc}
\usepackage[T1]{fontenc}
\usepackage{hyperref}
\title{Kott Language Resources for K-12 Curriculum Development}
\author{Compiled by Claude Sonnet 5.5 (claude-sonnet-5-5) for Andr\'as Kornai}
\date{}
\begin{document}
\maketitle

\section{Introduction}
Kott (ISO 639-3 zko) is an extinct Yeniseian language once spoken along the Mana River in central Siberia. It died out in the 1850s. The main historical source is M.~A.~Castr\'en's 1858 grammar. No school materials, revitalization programmes or audio were found in the searches, so only a scholarly base exists.

\section{Primary Educational Resources}
No primers, textbooks or school programmes were found.

\section{Linguistic and Reference Materials}
\begin{itemize}
\item Wikipedia article on the Kott language: \url{https://en.wikipedia.org/wiki/Kott_language}
\item FreeLang blog post, ``Kott: An extinct Central Siberian language'': \url{https://www.freelang.net/blog/?p=718}
\item Paper on kinship terminology (PDF, University of Naples ``L'Orientale''): \url{https://unora.unior.it/retrieve/dfd1bedd-b1a9-d55a-e053-3705fe0af723/kinship%20terminology.pdf}
\item Academia.edu item found in the search: \url{https://www.academia.edu/40222676}
\end{itemize}
Print sources named in the search results are Castr\'en (1858), \emph{Versuch einer jenissei-ostjakischen und kottischen Sprachlehre}; Werner (1997), \emph{Abri\ss{} der kottischen Grammatik}; and Verner (1990), \emph{Kottskij jazyk}. Online copies were not confirmed.

\section{Religious and Community Resources}
No religious or community resources were found.

\section{Digital and Audio Resources}
No audio, video or software resources were found.

\section{Curriculum Development Context}
Kott is documented only in a small body of 19th-century material. A heritage unit could draw on its place in the Yeniseian family, to which Ket belongs. A language-medium curriculum is not feasible with the evidence found.
\end{document}